where, if @@M0@@ is a vector space over @@M1@@, we write @@M2@@ for @@M3@@. For @@M4@@, @@M5@@ is surjective by (ii); therefore, if @@M6@@ is surjective, so is @@M7@@. By Serre's theorem, @@M8@@ is surjective for @@M9@@; hence (iii).

\sgaStableTarget{sga6.e13.prop.1-4}\textbf{Proposition 1.4.} \emph{Let @@M10@@ be an exact sequence on @@M11@@. If @@M12@@ is @@M13@@-regular, then we have}~:
\begin{enumerate}[label=(\roman*)]
\item \emph{@@M14@@ for @@M15@@ and @@M16@@.}
\item \emph{@@M17@@ for @@M18@@.}
\item \emph{@@M19@@ for @@M20@@.}
\end{enumerate}
\emph{In particular, @@M21@@ is @@M22@@-regular.}

Indeed, in the exact sequence
@@M23@@
the two outer groups are @@M24@@ for @@M25@@ and @@M26@@ by 1.3 (i). Thus @@M27@@, but, by Serre's theorem, @@M28@@ for @@M29@@; hence (i).

For @@M30@@, consider the following exact sequence:
@@M31@@
It implies (ii). Now suppose that @@M32@@ is surjective and consider the following commutative diagram:
@@M33@@
